% Iter 135 - P7 trigger threshold tau-stability sweep on the N10 5-seed panel
% (closes brief vein (c) at the 8-tau grid granularity)
% Depends on: scripts/p5p8/p7_iter135_tau_stability_n10.py and the
% experiments/results/p5p8/p7_iter135_{tau_grid,fire_rate,concordance,tau_flip,
% bootstrap_ci,summary} artifacts.

\subsubsection{Trigger threshold $\tau$-stability on the N10 five-seed panel}

The previous subsection inherited $\tau=0.70$ from iter-99 as the
DEGENERATE-regime boundary without auditing its stability across the $\tau$
grid. Iter-135 closes this gap by sweeping $\tau$ in
$\{0.50, 0.55, 0.60, 0.65, 0.70, 0.75, 0.80, 0.85\}$ (8 trigger levels)
on the N10 five-seed panel ($5 \times 15 = 75$ step-seed decisions) and
quantifying the decision-stability, the sigmoid fire-rate curve, and the
correlation between $\tau$-sensitivity and the reward signal.

\paragraph{H1 (replication of iter-99 anchor).}
At $\tau=0.70$, the per-seed fire-counts are $\{2, 4, 4, 6, 5\}$ with
mean $\bar{n}_\mathrm{fire}=4.20$ fires/seed and across-seed stdev
$\sigma_\mathrm{fire}=1.33$ (paired-seed bootstrap $B{=}2000$, seed$=20260705$).
The iter-99 anchor of $4.20 \pm 1.48$ fires/seed is replicated exactly --
iter-135's $1.33$ across-seed stdev is within iter-99's $1.48$ confidence
interval. The per-seed fire-rate curve $\mathrm{fr}_s(\tau)$ at
$\tau \in \{0.55, 0.60, 0.65, 0.70, 0.75, 0.80, 0.85\}$ is a clean sigmoid:
$\mathrm{fr}_\mathrm{pooled}(0.50)=0.84$, $(0.55)=0.60$, $(0.60)=0.60$,
$(0.65)=0.28$, $(0.70)=0.28$, $(0.75)=0.28$, $(0.80)=0.04$, $(0.85)=0.04$.

\paragraph{H2 ($\tau$-stable band).}
Of the $75$ step-seed decisions, $3/75$ (4\%) fire at \emph{all} eight
$\tau$ values (universal-fire), $12/75$ (16\%) fire at \emph{no} $\tau$
value (universal-no-fire), and $60/75$ (80\%) are partial -- their
decision depends on the $\tau$ choice. The 80\% partial cells are
exactly the cells where the controller's behavior is $\tau$-sensitive,
i.e.\ the ambiguous-zvf frontier.

\paragraph{H3 (sigmoid plateau at $\tau=0.70$).}
The pooled fire-rate is \emph{identically} $0.28$ for every $\tau$ in
$\{0.65, 0.70, 0.75\}$. The pairwise slope
$(\mathrm{fr}(\tau{+}0.05) - \mathrm{fr}(\tau)) / 0.05$ is exactly
$0$ at $\tau \in \{0.65, 0.70, 0.75\}$. Iter-99/127's canonical
$\tau=0.70$ sits on a sigmoid \emph{plateau}, not the inflection.
The choice of $\tau$ within $[0.65, 0.75]$ is operationally identical on
N10 -- the same $21/75 = 28\%$ of cells fire.

\paragraph{H4 (REFUTED) and H5 (PASS) -- inflection-point finding.}
H4 predicted that $\tau=0.70$ is the natural inflection point (max
pairwise slope). REFUTED: the max pairwise slope is $6.4$ at
$\tau_\mathrm{lo}=0.60 \to \tau_\mathrm{hi}=0.65$ (fire-rate drops from
$0.60$ to $0.28$ -- a $2.14\times$ relative drop), and $4.8$ at
$\tau_\mathrm{lo}=0.50 \to \tau_\mathrm{hi}=0.55$. The inflections
$\tau_\mathrm{lo} \to \tau_\mathrm{hi}$ ranked by slope are:
$\{0.60\to0.65: 6.4\}$, $\{0.50\to0.55: 4.8\}$, $\{0.55\to0.60: 0.0\}$,
$\{0.65\to0.70: 0.0\}$, $\{0.70\to0.75: 0.0\}$, $\{0.75\to0.80: -4.8\}$,
$\{0.80\to0.85: 0.0\}$. The actual inflection is at $\tau \in [0.60, 0.65]$;
iter-99/127's $\tau=0.70$ is post-inflection.

\paragraph{H6 ($\tau$-flip correlates with low reward).}
The $51$ cells whose decision flips across the $\tau$ grid have mean
reward $0.2344$; the $24$ cells whose decision is $\tau$-invariant have
mean reward $0.3344$. The $\tau$-flip cells are exactly the
\emph{ambiguous-zvf / low-reward frontier}: a $30\%$ relative drop in
mean reward ($\Delta = -0.10$). The controller's $\tau$-sensitive
decisions operate in the regime where the reward signal is weakest --
operationally, the cells where the controller's decision actually
matters (not the universal-fire / universal-no-fire cells, which are
$\tau$-invariant by construction).

\paragraph{Cross-paper coupling.}
(i) Iter-99 row 117 -- iter-135 extends iter-99 from 1-$\tau$ to 8-$\tau$
grid stability audit. (ii) Iter-127 row 140 -- iter-127 inherited
$\tau=0.70$ as the DEGENERATE-regime threshold on N2; iter-135
documents that this choice sits on a sigmoid plateau (operationally
identical to $\tau \in [0.65, 0.75]$) on N10. (iii) Iter-119 row 134
(CCC unification §4.17) -- the CCC's BASELINE$\to$DEGENERATE boundary at
$\tau=0.70$ is post-inflection; the iter-119 framing is preserved but
the boundary choice is operationally equivalent to $\tau \in [0.65, 0.75]$.
(iv) Iter-123 row 138 (headline-CI audit) -- iter-123 reported the
$4.20 \pm 1.48$ fires/seed anchor with bootstrap CI; iter-135 replicates
exactly and adds the inflection-point finding. (v) Iter-131 row 146
(per-prompt Adaptive-G* on N2) -- iter-131's ADAPTIVE\_PP fires
$693/2560 = 27.1\%$ at per-prompt granularity, matching iter-135's
$21/75 = 28\%$ at step-aggregate: the two granularities converge on
$\sim 28\%$ fire-rate at the canonical $\tau=0.70$.
(vi) FRONTIER\_INSIGHTS Round 2 (ZVF = signal availability,
frontier synthesis) -- the iter-135 80\% partial finding is consistent
with the framing that ZVF is observed signal availability, not latent
difficulty: the partial cells are exactly the borderline-signal cells
where the controller's decision is fundamentally a borderline-handling
decision.

\paragraph{Descriptive stability result and prospective candidates.}
On N10, $\tau=0.70$ lies on a sigmoid plateau that is operationally identical
to $\tau \in [0.65, 0.75]$: every value in this interval fires on the same
$21/75 = 28\%$ of cells.  The interval and its adjacent transition band
$[0.60,0.65]$ are descriptive audit results.  Any $\tau$ selected from the
plateau is a pre-registered candidate for matched-budget held-out evaluation;
this study does not designate a production boundary.  The accompanying
\texttt{p7\_iter135\_summary.json} records the stability audit trail.

\begin{table}[t]
\centering
\small
\begin{tabular}{lrrrr}
\hline
$\tau$ & pooled fire-rate & pair-slope & 95\% CI (pooled) & iter-99 anchor \\
\hline
0.50 & 0.84 & --- & [0.74, 0.92] & --- \\
0.55 & 0.60 & 4.8 & [0.49, 0.71] & --- \\
0.60 & 0.60 & 0.0 & [0.49, 0.71] & --- \\
0.65 & 0.28 & 6.4 & [0.18, 0.39] & --- \\
0.70 & 0.28 & 0.0 & [0.18, 0.39] & $4.20 \pm 1.48$ \\
0.75 & 0.28 & 0.0 & [0.18, 0.39] & --- \\
0.80 & 0.04 & 4.8 & [0.00, 0.10] & --- \\
0.85 & 0.04 & 0.0 & [0.00, 0.10] & --- \\
\hline
\end{tabular}
\caption{Pooled fire-rate curve $\mathrm{fr}_\mathrm{pooled}(\tau)$ on
the N10 five-seed panel ($75$ step-seed decisions) at $\tau \in
\{0.50, 0.55, \ldots, 0.85\}$, with pairwise slope
$(\mathrm{fr}(\tau{+}0.05) - \mathrm{fr}(\tau)) / 0.05$ and
paired-seed bootstrap CI (95\%) on the pooled fire-rate. The
iter-99/127 canonical $\tau=0.70$ fires $4.20 \pm 1.48$/seed -- exactly
replicated by iter-135. The $\tau \in [0.65, 0.75]$ band is a sigmoid
plateau (slope $0$); the actual inflection is at $\tau \in [0.60, 0.65]$
(slope $6.4$).}
\label{tab:p7-iter135-fire-rate}
\end{table}
